\begin{enumerate}
\item We know $V_1=\omega a_1$ and $V_2=\omega a_2$. \hfill(1 point)
From Kepler's third law, $P^2=4\pi^2a^3/(GM)$, where $M=m_1+m_2$ and $a=a_1+a_2$. \hfill(2 points)
Also $P=2\pi/\omega$, $a_1=m_2a/(m_1+m_2)$ and $a_2=m_1a/(m_1+m_2)$. \hfill(1 point)
Replacing these values in Kepler's equation gives
\[G(m_1+m_2)=4\pi^2\frac{\omega^2(m_1+m_2)^3}{4\pi^2m_2^3}a_1^3,\]
\[\frac{m_2^3}{(m_1+m_2)^2}=(\omega a_1)^3\frac P{2\pi G}.\]
\hfill(2 points)
When we observe a star we measure the line-of-sight component, rather than the real orbital velocity. \hfill(2 points)
The required mass functions are
\[\frac{m_2^3}{(m_1+m_2)^2}\sin^3\theta=v_{1,\mathrm{obs}}^3\frac P{2\pi G},\]
\hfill(1 point)
\[\frac{m_1^3}{(m_1+m_2)^2}\sin^3\theta=v_{2,\mathrm{obs}}^3\frac P{2\pi G}.\]
\hfill(1 point)
Here $v_{1,\mathrm{obs}}$ is the maximum observed orbital velocity.
\item We know $v_1/v_2=m_2/m_1$. \hfill(1 point)
From the graph $v_1/v_2=1.5$, so $m_2=1.5m_1$. The period is 40 days and $v_{1,\mathrm{obs}}=75\,\mathrm{km\,s^{-1}}$. For the minimum value take $\sin\theta=1$ and put these values in the mass function. \hfill(2 points)
This gives $0.54m_1=1.74M_\odot$, hence $m_1=3.2M_\odot$ \hfill(1 point)
and similarly $m_2=4.9M_\odot$. \hfill(1 point)
The values are lower limits, since $\sin\theta\leq1$.
\end{enumerate}
